\section{Introduction}

\subsection{Contexte}
La carte nationale d'identité électronique est le titre régalien le plus diffusé du Royaume~: la génération lancée en août 2020 équipait déjà quinze millions de citoyens en 2023, soit environ 85~\% de la population de plus de seize ans \cite{dgsn2023tcn}, et la DGSN a produit près de 3,7 millions de cartes au cours de la seule année 2025 \cite{lavieeco2025bilan}. Au-delà du support physique, la loi 04-20 a fait de la CNIE un titre certifiant l'identité numérique de son porteur \cite{loi0420}. Sur cette base, la DGSN et l'Agence de Développement du Digital ont ouvert en avril 2022 la plateforme «~Tiers de Confiance National~», qui permet à des fournisseurs de services accrédités de déléguer l'identification et l'authentification de leurs usagers \cite{lavieeco2022tcn}. Plus de trente organismes publics et privés y étaient raccordés fin 2024 \cite{allafrica2024bilan}, et le portail \texttt{epolice.ma}, lancé en décembre 2024, s'appuie sur elle pour délivrer la fiche anthropométrique à distance \cite{telquel2024epolice}.

Le portail \texttt{cnie.ma}, objet de cette étude, constitue le point d'entrée du processus d'obtention et de renouvellement de la carte~: depuis septembre 2020, la pré-demande et la prise de rendez-vous en ligne y sont obligatoires avant tout passage en centre d'enregistrement \cite{telquel2020rdv}. Ce dispositif a absorbé une forte montée en charge, mais il reste un portail de «~première génération~»~: il collecte des données puis renvoie le citoyen vers un parcours physique largement manuel.

\subsection{Problématique et commande}
Les décideurs de la DGSN souhaitent faire évoluer ce système pour le rendre plus efficace, plus réactif et plus personnalisé, en y intégrant des outils d'intelligence artificielle. Dans le rôle du cabinet d'urbanisation mandaté, nous devons~: (i)~étudier et cartographier l'existant, (ii)~proposer le système cible en nous appuyant sur l'existant, sur les standards d'architecture d'entreprise et sur les bonnes pratiques du domaine de l'identité numérique.

La question centrale est donc la suivante~: \emph{comment réorganiser le système d'information du périmètre CNIE pour que le parcours du citoyen devienne majoritairement numérique, que la fraude soit détectée au plus tôt, que la capacité des centres soit pilotée par la donnée, tout en respectant les exigences de souveraineté et de protection des données personnelles~?}

\subsection{Périmètre et choix structurants}
Le processus principal retenu est \textbf{P0 -- Demande et renouvellement de la CNIE}, de l'expression du besoin jusqu'à la remise de la carte et l'activation de l'identité numérique. Il concentre l'essentiel des volumes, des irritants et des opportunités d'automatisation. Le sous-processus \textbf{SP1 -- Pré-demande en ligne, contrôle des pièces et prise de rendez-vous} a été choisi pour le prototypage sous Bonitasoft, car il est celui où l'intelligence artificielle transforme le plus directement l'expérience du citoyen.

\subsection{Organisation du rapport}
La section~\ref{sec:methodo} présente la démarche d'urbanisation et les outils employés. La section~\ref{sec:existant} cartographie l'existant et en dresse le diagnostic. La section~\ref{sec:strategie} capte la stratégie et les cas d'usage d'intelligence artificielle. Les sections~\ref{sec:metier} à~\ref{sec:technique} déclinent les cartographies métier, fonctionnelle, applicative (avec le prototype Bonita) et technique du système cible. La section~\ref{sec:trajectoire} propose la trajectoire de migration et la gouvernance, avant la conclusion.
